%%=============================================================================
%% Samenvatting
%%=============================================================================

%%---------- Nederlandse samenvatting -----------------------------------------
%
\IfLanguageName{english}{%
\selectlanguage{dutch}
\chapter*{Samenvatting}
\selectlanguage{english}
}{}

%%---------- Samenvatting -----------------------------------------------------
% De samenvatting in de hoofdtaal van het document

\chapter*{\IfLanguageName{dutch}{Samenvatting}{Abstract}}

Capture The Flag (CTF)-events zijn cybersecuritywedstrijden waarbij deelnemers technische uitdagingen oplossen om punten te behalen. Een deel van die uitdagingen vereist een netwerktoegankelijke service, bijvoorbeeld een webapplicatie of kwetsbare containeromgeving. Op on-premise infrastructuur worden dergelijke challenge-omgevingen vaak nog manueel of met losse scripts beheerd. Dat verhoogt de operationele belasting voor organisatoren en vergroot het risico op misconfiguraties, onvoldoende isolatie tussen deelnemers en resource-uitputting op de host.

Deze bachelorproef onderzoekt hoe Docker-containers die gekoppeld zijn aan CTFd-challenges dynamisch kunnen worden beheerd, aangemaakt en beëindigd op basis van gebruikersacties en configureerbare parameters binnen een volledig self-hosted context. De centrale onderzoeksvraag luidt hoe dit veilig, reproduceerbaar en schaalbaar kan gebeuren zonder afhankelijkheid van betaalde cloudcomponenten of Kubernetes als vereiste.

De methodologie combineert een literatuurstudie met praktijkgericht en iteratief onderzoek. Eerst werden CTFd, Docker, de Docker Engine API en direct toepasbare containerbeveiligingsmaatregelen bestudeerd. Daarna werd een verkennende standalone proof-of-concept gebouwd om container lifecycle management, timeout cleanup, netwerkisolatie, resource-limieten en idempotent startgedrag los van CTFd te valideren. Op basis van die bevindingen werd vervolgens een eigen CTFd-plugin ontwikkeld die een nieuw challenge-type toevoegt en de containerlevenscyclus rechtstreeks koppelt aan acties van spelers en beheerders.

De uiteindelijke plugin-PoC ondersteunt challengeconfiguratie met image, een of meerdere containerpoorten, CPU- en geheugenlimieten, timeout en een maximum aantal gelijktijdige instanties. Daarnaast ondersteunt de plugin het uploaden van Docker-archieven, het automatisch herimporteren van die archieven wanneer een image lokaal ontbreekt, het hergebruiken van een bestaande instantie voor dezelfde speler of hetzelfde team, automatische cleanup bij solve of timeout en een administratieve runtime-pagina binnen CTFd. Voor elke runtime worden basismaatregelen toegepast zoals een read-only root filesystem, capability dropping, \texttt{no-new-privileges}, PID-limieten, CPU- en geheugenlimieten en een afzonderlijk bridge-netwerk per instantie.

De validatie gebeurde in meerdere stappen. De standalone PoC werd gevalideerd met \texttt{pytest}, een Docker-gebaseerde smoke test, isolatie-experimenten, timeoutcontrole en een load snapshot met Locust. De plugin-PoC werd daarna end-to-end getest in same-host modus, in een lokale split-host simulatie en in een reële opstelling met twee afzonderlijke Lima-VM's. Die validatie toont aan dat een CTFd-plugin zowel een lokale Docker-daemon als een remote Docker-host kan aansturen, terwijl spelers correcte runtime-URL's ontvangen voor de afzonderlijke challenge-host.

De conclusie van dit onderzoek is dat een plugin-gebaseerde aanpak binnen CTFd een haalbare en bruikbare oplossingsrichting vormt voor dynamisch containerbeheer bij CTF-events op eigen infrastructuur. Tegelijk blijkt dat verdere productierijpe uitwerking nog aandacht vraagt voor TLS-beveiligde remote Docker-communicatie, uitgebreidere foutinjectie, reverse-proxy-gebaseerde routing en schaalverdeling over meerdere runtime-hosts.
